% Integración por delta de la adenda de realizaciones físicas.
% Residencia de procedencia: tipo dimensional de G, selector de holonomía y carta decimal.

\chapter{Transduction dimensionnelle et sélection holonomique de la constante
gravitationnelle}
\label{chap:transductor-gravitatorio-G}
\index{constante gravitationnelle!transducteur dimensionnel}
\index{naturalité!sélecteur gravitationnel}
\index{lecteur décimal!contrôle de G}

\section{Droite gravitationnelle et forme dimensionnelle}
\label{sec:recta-gravitatoria-interna}

Dans le réseau du \cref{chap:funtor-dimensional-hmt}, la dimension de la
constante gravitationnelle est
\[
 \mathbf G=-\mathbf S+5\mathbf C+2\mathbf T,
 \qquad
 [G]=L^3M^{-1}T^{-2}.
\]
Soient
\[
 c_{\rm int}\in\mathcal L_C^+,\qquad
 t_0\in\mathcal L_T^+,\qquad
 \ell_0:=c_{\rm int}t_0\in\mathcal L_L^+,
 \qquad
 \hbar_{\rm int}\in\mathcal L_S^+.
\]
La section \(\hbar_{\rm int}\) n'est obtenue qu'après le lecteur
déterminantal de décade \(-34\) et le choix d'une base de la droite
d'action. Ce n'est pas le caractère angulaire adimensionnel antérieur au retour.

Pour tout caractère adimensionnel \(\theta_G>0\), nous définissons
\[
 \boxed{
 G_{\rm int}(\theta_G)
 =\theta_G
 \frac{c_{\rm int}^{5}t_0^2}{\hbar_{\rm int}}
 =\theta_G
 \frac{c_{\rm int}^{3}\ell_0^2}{\hbar_{\rm int}}
 \in\mathcal L_G^+.}
\]
L'égalité des deux expressions et leur type sont
\textsc{exacts internes} dans le réseau dimensionnel. Ils ne sélectionnent pas
\(\theta_G\). Si
\[
 L_G^2:=\frac{\hbar_{\rm int}G_{\rm int}}{c_{\rm int}^3},
\]
alors
\[
 \theta_G=\left(\frac{L_G}{\ell_0}\right)^2.
\]
Cette dernière identité localise la donnée de sélection : le rapport des longueurs doit
être produit par une structure d'holonomie ou de réponse spectrale
antérieure à la carte métrologique.

\section{Obstruction du caractère positif sur le retour observable}
\label{sec:selector-holonomia-G}

Le \cref{sec:jerarquia-retornos-27-54-108} définit
\[
 \mathcal K_{27}:=F_{\rm obs}^{27}.
\]
La proposition primaire \cref{prop:retorno-fobs-identidad} démontre
\[
 F_{\rm obs}^{108}=I;
\]
par conséquent, dans le quotient observable,
\[
 \mathcal K_{27}^{4}=I.
\]

\begin{theorem}[Obstruction de torsion pour le sélecteur positif observable]
\label{thm:no-go-chi-G-observable}
Tout homomorphisme multiplicatif
\[
 \chi:\langle\mathcal K_{27}\rangle
 \longrightarrow\mathbb R_{>0}^{\times}
\]
est trivial. En particulier,
\[
 \chi(\mathcal K_{27})=1.
\]
\end{theorem}

\begin{proof}
La multiplicativité et \(\mathcal K_{27}^{4}=I\) impliquent
\[
 \chi(\mathcal K_{27})^4=\chi(I)=1.
\]
Le seul élément réel positif dont la puissance quatrième est un est \(1\).
\end{proof}

L'ordre observable est exactement quatre, et non simplement un diviseur de quatre :
dans la coordonnée \(\theta\in\mathbb Z/108\mathbb Z\),
\[
 \mathcal K_{27}(\theta)=\theta+27,
 \qquad
 \mathcal K_{27}^{2}(\theta)=\theta+54\neq\theta,
 \qquad
 \mathcal K_{27}^{4}(\theta)=\theta.
\]
Le même groupe cyclique admet des caractères unitaires
\[
 \chi_k(\mathcal K_{27})
 =\exp\!\left(\frac{2\pi i k}{4}\right),
 \qquad k=0,1,2,3,
\]
et, en particulier, les valeurs de phase \(\pm i\). La conclusion typée est
que le cycle observable peut transporter une phase, mais non une échelle positive
non triviale. Ces caractères de phase ne s'identifient pas à l'opérateur
\(J_W\) de Paley : ils satisfont une relation quadratique analogue sur des domaines
distincts.

L'obstruction ne dépend pas d'une comparaison métrologique. Elle affecte la forme
même du sélecteur : un caractère positif sur le cycle observable ne peut
produire un \(\theta_G\) non trivial. En outre, un simple ensemble de classes de
conjugaison ne constitue pas, en général, un groupe et ne peut être déclaré domaine
d'un caractère sans que sa composition soit définie au préalable.

\subsection{Unique continuation multiplicative admissible}

Une voie multiplicative ne peut être rouverte que sur un objet enrichi qui
n'hérite pas de la torsion observable. Pour un objet \(x\) du groupoïde de mémoire,
le domaine correctement typé serait le groupe d'isotropie
\[
 H_x:=\operatorname{Aut}_{\mathcal G_{\rm mem}}(x)
\]
et, par multiplicativité vers un groupe abélien, son abélianisation
\[
 H_x^{\rm ab}\longrightarrow\mathbb R_{>0}^{\times}.
\]
Tout caractère positif annule la partie de torsion de \(H_x^{\rm ab}\). Ainsi,
la proposition doit construire un relèvement
\[
 \widetilde{\mathcal K}_{27}\in H_x,
 \qquad
 p(\widetilde{\mathcal K}_{27})=\mathcal K_{27},
\]
prouver que sa classe abélienne est non de torsion et sélectionner naturellement
un caractère \(\widetilde\chi_G\). La voie positive sur le cycle observable
est ainsi fermée par obstruction : sans cette donnée non de torsion, il n'existe pas de
sélecteur multiplicatif non trivial. Si tout relèvement admissible
satisfaisait
\(\widetilde{\mathcal K}_{27}^{\,4}=I\), la voie multiplicative serait
également écartée dans le domaine enrichi.

\section{Voie chronologique de longueur, d'action et de spectre}
\label{sec:via-cronologica-espectral-G}

Le \cref{thm:cobertura-acumulativa-dirigida} fournit une troisième voie,
distincte tant du caractère observable obstrué que de l'éventuel
relèvement dans un groupe d'isotropie. Sur
\[
 \widetilde{\mathcal O}_{108}^{+}
 =\mathcal O_{108}\times\mathbb N^2,
\]
le retour de vingt-sept pas ajoute \((1,18)\) et sa puissance quatrième ajoute
\((4,72)\). Cette action appartient à une catégorie chronologique dirigée. Ce n'est pas
un relèvement du groupoïde ni un champ dont il aurait été démontré qu'il est natif de
l'état TPK enrichi complet.

La fonctionnelle positive déjà construite,
\[
 \mathcal A_\ast(\gamma)
 =G_{\rm or}(\gamma)+\lambda_\ast S(\gamma),
\]
est additive par concaténation d'histoires dirigées. Ce n'est pas un caractère du
groupe d'isotropie, car ses indicateurs ne changent pas de signe sous une
inversion formelle. La complétion
\(\mathbb N^2\to\mathbb Z^2\) permet d'introduire des caractères unitaires pour
l'analyse harmonique, mais n'altère pas ce fait.

\begin{equation}
\label{eq:fraccion-angular-orientada-G}
 \delta_\theta
 :=\frac{A_{\rm deg}-C_{\rm deg}^{\ast}}{360}.
\end{equation}
Cette fraction angulaire orientée est définie une fois construits l'angle et le
contre-angle uniques. Elle sera utilisée dans les deux gabarits dimensionnels
ultérieurs et n'intervient pas rétroactivement dans l'obtention de
\(A_{\rm deg}\) ni de \(C_{\rm deg}^{\ast}\).

\subsection{Opérateur tordu et contrôle négatif déterministe}

Pour le contrôle exact, soit \(\Gamma=(V,E)\) un multigraphe orienté
\emph{fini} enrichi, muni d'une mesure
\(\mu:V\to\mathbb R_{>0}\), de poids \(a_\ast(e)\geq0\), d'un paramètre
\(\beta\geq0\) et d'un cocycle
\(c(e)\in\mathbb Z^2\). Pour une normalisation fixée, considérons
\[
 (\mathcal L_{\beta,\vartheta}f)(y)
 =
 \sum_{e:x\to y}
 e^{-\beta a_\ast(e)}
 e^{i\langle\vartheta,c(e)\rangle}f(x),
 \qquad
 \vartheta\in\mathbb T^2,
\]
et prenons pour domaine tout l'espace de Hilbert fini
\(\ell^2(V,\mu)\). Alors \(\mathcal L_{\beta,\vartheta}\) est borné et
\[
 \mathcal H_{\beta,\vartheta}
 :=\mathcal L_{\beta,\vartheta}^{*}
   \mathcal L_{\beta,\vartheta}\geq0
\]
est autoadjoint. Une extension à un graphe infini exigerait, en outre,
des opérateurs bornés ou bien des opérateurs fermés densément définis tels
que \(\mathcal L_{\beta,\vartheta}^{*}
\mathcal L_{\beta,\vartheta}\) soit autoadjoint sur un domaine commun ;
cette extension n'est pas présupposée ici.

\begin{proposition}[Annulation de phase dans l'orbite déterministe]
\label{prop:control-negativo-espectral-determinista}
Si \(\Gamma\) est uniquement l'orbite déterministe observable et que chaque sommet
possède une seule arête entrante, alors
\(\mathcal H_{\beta,\vartheta}\) est indépendant de \(\vartheta\). En
particulier, la hessienne de chacune de ses bandes par rapport à
\(\vartheta\) est nulle.
\end{proposition}

\begin{proof}
Dans chaque composante, \(\mathcal L_{\beta,\vartheta}\) multiplie l'unique
terme entrant par une phase unitaire. Dans
\(\mathcal L_{\beta,\vartheta}^{*}\mathcal L_{\beta,\vartheta}\), cette phase
est multipliée par sa conjuguée et s'annule. Il ne reste pas de termes croisés
entre histoires distinctes.
\end{proof}

Ce contrôle empêche d'attribuer une rigidité spectrale au simple cycle déterministe.
Une réponse de phase non triviale exige au moins l'une des constructions
suivantes :
\begin{enumerate}
 \item un multigraphe enrichi dans lequel deux histoires cohérentes ou plus
 contribuent à une même amplitude ; ou
 \item une différentielle magnétique
 \[
  (d_\vartheta f)(e)
  =
  e^{i\langle\vartheta,c(e)\rangle/2}f(t(e))
  -
  e^{-i\langle\vartheta,c(e)\rangle/2}f(s(e)),
 \]
 avec un laplacien
 \(\Delta_\vartheta=d_\vartheta^*d_\vartheta\), dont la phase dépende de
 l'holonomie de cycles enrichis.
\end{enumerate}
Dans les deux cas, l'espace de Hilbert, le domaine, la mesure,
les poids et la normalisation doivent être déclarés, et il faut démontrer que la famille n'est pas
unitairement équivalente à la famille sans phase.

Supposons que
\(\vartheta\mapsto\mathcal H_{\beta,\vartheta}\) soit une famille
autoadjointe analytique sur un domaine commun et que
\(\lambda_0(\vartheta)\) soit une bande simple et isolée dans un voisinage de
\(\vartheta=0\). Si \(0\) est un point critique et un minimum local,
\[
 Q_{ij}
 =
 \left.
 \frac{\partial^2\lambda_0(\vartheta)}
      {\partial\vartheta_i\partial\vartheta_j}
 \right|_{\vartheta=0}
\]
est semi-définie positive. Sans l'hypothèse de minimum, \(Q\) représente seulement
la réponse quadratique. Pour
\[
 v_{27}=(1,18),
 \qquad
 \mathcal R_{27}=v_{27}^{\mathsf T}Qv_{27},
\]
la condition \(\mathcal R_{27}\neq0\) est une preuve minimale de sensibilité
spectrale à la mémoire accumulée.

\subsection{Normalisation typée de la réponse spectrale et construction de \texorpdfstring{\(\theta_G^{\mathrm{spec}}\)}{theta G spec}}

Si une normalisation naturelle et unique \(\mathfrak N\) transforme la réponse
précédente en un rapport positif indépendant de l'échelle globale de
l'opérateur, on peut définir
\[
 \Lambda_G
 :=\ell_0\,\mathfrak N(Q,v_{27}),
 \qquad
 \theta_G^{\rm spec}
 :=
 \left(
 \delta_\theta\frac{\Lambda_G}{\ell_0}
 \right)^2,
\]
et, par le transducteur dimensionnel,
\[
 \boxed{
 G_{\rm spec}
 =
 \frac{c_{\rm int}^{3}}{\hbar_{\rm int}}
 \bigl(\delta_\theta\Lambda_G\bigr)^2.}
\]
L'homogénéité du gabarit est exacte. La classe de ses évaluations est
caractérisée par les conditions de sélection canonique de l'opérateur,
de réponse non nulle, de naturalité et de stabilité sous raffinement, d'indépendance
vis-à-vis du changement d'échelle, d'unicité de \(\mathfrak N\) et de scellement préalable à toute consultation
de la valeur métrologique de \(G\). Sans ces conditions, l'invariance par
changement d'échelle démontre la non-unicité du sélecteur ; avec elles, l'évaluation est
déterminée dans le domaine normalisé déclaré.

\section{Fonctionnelle non homomorphe sur les voies enrichies}
\label{sec:funcional-longitud-K27-G}

La seconde continuation possible abandonne le caractère multiplicatif et utilise
une fonctionnelle géométrique, spectrale ou d'action. Soit
\[
 \mathscr L:
 \operatorname{Loop}(\mathcal X_{\rm TPK}^{\rm enr})
 \longrightarrow\mathcal L_L^+\cup\{0\}
\]
invariante par changement cyclique du point de base, réétiquetage admissible,
conjugaison et subdivision. Sa loi de concaténation doit être déclarée
explicitement ---par exemple, comme norme sous-additive ou fonctionnelle spectrale--- ;
on ne présume pas qu'elle soit un homomorphisme du groupe de torsion. Pour un
relèvement enrichi candidat, nous écrivons
\[
 \Lambda_{27}:=\mathscr L(\widetilde{\mathcal K}_{27}),
 \qquad
 \rho_{27}:=\frac{\Lambda_{27}}{c_{\rm int}t_0}.
\]
La fraction angulaire orientée de
\eqref{eq:fraccion-angular-orientada-G} produit le candidat adimensionnel
\[
 \theta_G^{\rm hol}:=(\delta_\theta\rho_{27})^2.
\]
Sa spécialisation dans le transducteur donne
\[
 \boxed{
 G_{\rm hol}
 =\frac{c_{\rm int}^{3}}{\hbar_{\rm int}}
  \bigl(\delta_\theta\Lambda_{27}\bigr)^2.}
\]
La forme temporelle est la même identité :
\[
 T_{27}:=\frac{\Lambda_{27}}{c_{\rm int}},
 \qquad
 t_\ast:=\delta_\theta T_{27}.
\]
En particulier, il ne convient pas de remplacer \(t_\ast\) par
\(T_{27}/\delta_\theta\).

L'homogénéité de ces formules est exacte et cette carte conserve le statut
de réalisation fonctionnelle auxiliaire. Le propriétaire gravitationnel complet
construit ensuite, à partir de la branche d'action déjà générée, le sélecteur unique
\(\vartheta_{108}=(54/\pi_{\mathrm{HMT}})^2\), la section \(G_a\) et
l'invariant typé \(G_a\hbar_a\). Cette construction gouverne la dérivation de
\(G\) ; la fonctionnelle \(\mathscr L\) enregistre une coordinatisation supplémentaire de
la même sortie. L'ordre causal conserve d'abord \(A\) et \(C^\ast\), puis
la transduction gravitationnelle et enfin la carte métrologique.

\section{Transduction de la carte décimale gravitationnelle sur la droite de grandeur et séparation causale de la reconnaissance métrologique}
\label{sec:lector-decimal-independiente-G}

Soit
\[
 \mathscr R_G([n,k,d]_3;t,e)
 :=10^{-n}\left(k+\frac{t}{10^3}+\frac{e}{10^d}\right),
\]
où l'indice de \([n,k,d]_3\) enregistre le triplet du lecteur et ne désigne pas
un numéral ternaire. Pour les données déclarées
\[
 [n,k,d]_3=[11,6,5]_3,\qquad t=674,\qquad e=30,
\]
l'évaluation rationnelle est
\[
 \boxed{
 \mathscr R_G
 =10^{-11}\left(6+\frac{674}{1000}+\frac{30}{10^5}\right)
 =6.67430\times10^{-11}.}
\]
L'égalité arithmétique est exacte une fois les paramètres déclarés. Le triplet,
les blocs \(674\), \(30\) et leur position décimale ont le statut de données
de la carte, de sorte que cette branche est conservée comme contrôle paramétré et
reconnaissance ultérieure. Le nombre est adimensionnel jusqu'au choix d'une base
\(u_G\in\mathcal L_G^+\) ; alors seulement la section est définie
\[
 G_{\rm dec}:=\mathscr R_G\,u_G.
\]
La coïncidence de sa coordonnée avec une publication métrologique est une
\textsc{reconnaissance externe}. Les entiers \(674\), \(30\), le triplet
\([11,6,5]\) et le placement décimal appartiennent exclusivement à cette carte
de reconnaissance ; la dérivation gouvernante réside dans le sélecteur
\(\vartheta_{108}\) et dans le transducteur d'action du propriétaire complet.

\section{Hiérarchie probatoire des constructions gravitationnelles}
\label{sec:cuadrado-contraste-G}

Les couches précédentes ne constituent pas deux généalogies prouvées d'une même
constante. Leur situation est la suivante :
\begin{center}
\begin{tabularx}{0.96\textwidth}{@{}L{0.25\textwidth}Y L{0.26\textwidth}@{}}
\toprule
Objet & Résultat disponible & Verdict\\
\midrule
Transductor dimensional &
Type exactement \(G_{\rm int}(\theta_G)\), sans fixer \(\theta_G\). &
Théorème dimensionnel ; non sélecteur.\\
Caractère de l'observable \(\mathcal K_{27}\) &
\(\mathcal K_{27}^4=I\) impose
\(\chi(\mathcal K_{27})=1\). &
Formulation obstruée.\\
Levantamiento enriquecido &
Possible seulement avec une classe abélienne non de torsion, un caractère naturel et une
normalisation unique. &
Obstruction de canonicité démontrée sans ces données.\\
Revêtement chronologique et réponse spectrale &
La translation cumulative est d'ordre de semi-groupe infini ; dans l'orbite
déterministe, la phase s'annule. &
Réponse nulle démontrée dans l'orbite déterministe ; classe non triviale
caractérisée par l'interférence et la normalisation naturelle.\\
Fonctionnelle de voie &
La formule de \(G_{\rm hol}\) est homogène une fois
\(\Lambda_{27}\) fixé. &
Famille paramétrée exacte ; il n'existe pas de sélecteur univoque sans spécifier
\(\mathscr L\) et sa normalisation.\\
Lecteur décimal &
Reproduit exactement la coordonnée déclarée, mais reçoit ses blocs et sa
position. &
Control parametrizado.\\
\bottomrule
\end{tabularx}
\end{center}

\begin{criterion}[Sélecteur de la constante gravitationnelle]
\label{crit:selector-G}
Une promotion est rejetée si :
\begin{enumerate}
 \item le relèvement enrichi conserve uniquement de la torsion ou ne descend pas
 à une classe abélienne bien définie ;
 \item l'opérateur chronologique ne conserve pas l'interférence, produit
 \(\mathcal R_{27}=0\), est dépourvu de domaine ou de mesure, ou sa réponse change sous
 un raffinement admissible ;
 \item \(\mathscr L\) dépend du point de base, du réétiquetage, de
 l'orientation choisie, d'une subdivision ou d'une normalisation libre ;
 \item la normalisation \(\mathfrak N\) dépend de l'échelle globale de
 l'opérateur, n'est pas unique ou est fixée en consultant \(G\) ;
 \item \(674\), \(30\) ou le placement décimal admettent des alternatives ayant le
 même droit structurel ;
 \item le candidat consulte la valeur métrologique avant d'être scellé ;
 \item la valeur n'est pas stable sous raffinement ou ne définit pas une section de
 \(\mathcal L_G\) dans une carte dimensionnelle commune.
\end{enumerate}
\end{criterion}

\section{Rapports gravitationnel--inertiel et constante dimensionnelle}
\label{sec:razones-no-rivales-G}

La normalisation locale de la carte torsionnelle,
\[
 \frac{\mathcal G_{\rm tor}}{\mathcal I_{\rm tor}}=1,
\]
et l'interface surfacique--volumétrique du
\cref{sec:ambivalencia-grav-inercial-pi-phi2},
\[
 \frac{\mathcal G_{\rm av}}{\mathcal I_{\rm av}}
 =\frac{\pi_{\rm HMT}}{\varphi^{2}},
\]
sont des rapports adimensionnels de lecteurs. Ce ne sont pas deux valeurs rivales de la
constante dimensionnelle \(G\). La première exprime une restriction diagonale
locale ; la seconde est une interface conditionnée par la factorisation à travers une
mémoire commune. La construction d'une section de
\(\mathcal L_G\) requiert, en outre, le repère
\((\hbar_{\rm int},c_{\rm int},\ell_0)\) et un sélecteur
\(\theta_G\).
